% frontmatter/verano.tex -- instrucciones para el adulto del cuaderno de
% verano (ediciones/main-verano.tex, ver notes/07-ediciones.md): van en
% vez de frontmatter/como-usar.tex, que habla del año entero. Detrás, en
% vez del mapa del curso, el mapa del verano (\mapaEdicion, preamble.tex).
\section*{Cómo usar este cuaderno de verano}

Este cuaderno es el verano de \emph{Aprendo a leer}, el cuaderno de
frases: sus últimas 13 semanas, en un cuaderno aparte, para leer un poco
cada día de las vacaciones. Está pensado para una niña o un niño que ya
lee frases completas y quiere ganar soltura -- no enseña a leer desde
cero. El verano es precisamente cuando más se agradece no perder el
ritmo de lectura.

Tiene 65 páginas, una para cada día de lunes a viernes, y todas siguen
la misma estructura:

\begin{itemize}[leftmargin=6mm]
\item \textbf{Fecha}, \textbf{Leído} y \textbf{Notas}: para el adulto.
  Marcar cada día si ha leído sola, con ayuda o con dificultad no cuesta
  nada, y al final del verano es un registro de todo lo que ha avanzado.
\item Las \textbf{cuatro frases del día}, dentro del recuadro verde:
  primero se leen en silencio y luego en voz alta. Algunos días hay
  diálogo, y algunos viernes (\textbf{\lblRelee}) está el texto de toda
  la semana junto, para leerlo otra vez de principio a fin.
\item Una \textbf{actividad} distinta cada día, siempre relacionada con
  lo que se acaba de leer:
  \begin{itemize}[leftmargin=6mm]
  \item \textbf{\lblResponde}: una pregunta corta sobre lo leído, para
    comprobar que se ha entendido y no solo que se ha leído.
  \item \textbf{\lblBusca}: encontrar un dato en el texto.
  \item \textbf{\lblRelaciona}: unir con una línea lo que va junto.
  \item \textbf{\lblAdivina}: una adivinanza sobre la historia de la
    semana.
  \item \textbf{\lblDibuja} y \textbf{\lblCompleta}: dibujar algo de lo
    que cuenta el texto, o convertir unas pocas líneas ya dibujadas en
    un dibujo -- no hay un dibujo ``correcto''.
  \item \textbf{\lblTraza}: repasar con el dedo o con un lápiz el
    contorno de una letra, mayúscula y minúscula, y luego escribirla
    sola en la línea.
  \item \textbf{\lblRepasa}: cada dos semanas, un pequeño resumen con
    casillas para marcar y una tarea final.
  \end{itemize}
\end{itemize}

Las páginas cuentan el verano de \textbf{Lucía}, su hermano
\textbf{Dani} y su perro \textbf{Toby} en el pueblo de la abuela
\textbf{Rosa}: el río, el huerto, una tormenta, la playa, la verbena, un
amigo nuevo -- y, al final, la vuelta a la ciudad y al cole. Cada semana
es un capítulo nuevo de la misma historia; en la página siguiente están
todas, con un sol para colorear cada día que se lee.

Cinco o diez minutos al día son suficientes, y no hace falta terminar
la actividad el mismo día -- lo importante es la lectura. Si se atasca
en casi cada palabra, que lea solo la primera frase y que el resto lo
lea un adulto en voz alta: lo importante es que la sesión siga siendo
agradable, no una lucha. Cada diez páginas hay un pequeño cartel de
celebración, y al final del cuaderno, una \textbf{clave de respuestas}
para las páginas de \lblAdivina\ y \lblRelaciona, y un diploma.

{\small\color{colorGris} \emph{Aprendo a leer} tiene también el año
entero, con sus 260 páginas, y otros dos niveles: \emph{Primeras
palabras}, el anterior, y \emph{Leo con lupa}, el siguiente.\par}
\newpage
